\section{The project graph}
\begin{frame}{Concrete first: a project with seven parts}
\princ{MODEL VS PROCEDURE}
\begin{columns}[T]
\begin{column}{0.52\textwidth}
Imagine a small software project. Its seven parts are:
\begin{itemize}
\item project, plan, build
\item requirements, design
\item tests, docs
\end{itemize}
These are synthetic labels, not a real project.
\end{column}
\begin{column}{0.43\textwidth}
\begin{block}{Question}
What must happen before what?
\end{block}
\predict{Which two parts can start directly after project?}
\end{column}
\end{columns}
\note{Keep the example concrete. Read the seven labels and say that they are invented for the lesson. Ask which items project points to directly. Students should predict plan and build. This is a dependency question, not yet a folder question.}
\end{frame}

\begin{frame}{A directed edge means dependency here}
\princ{EDGE MEANING}
\begin{block}{Our graph rule}
An arrow \texttt{A -> B} means: complete A before B.
\end{block}
\begin{itemize}
\item \texttt{project -> plan}: plan follows project.
\item \texttt{build -> tests}: tests follow build.
\item \texttt{requirements -> tests}: tests also depend on requirements.
\end{itemize}
\mistake{Treating every arrow as ``is stored inside'' changes the question.}
\see{the arrow has a direction and a work-order meaning.}
\note{Define the edge in plain language and give three examples. Pause on requirements to tests: it will cross branches later. Ask students to say the sentence with the arrow reversed; reversing it changes the claim. The common mistake is the exact model drift we are auditing.}
\end{frame}

\begin{frame}{The seven-edge dependency graph}
\princ{EDGE MEANING}
\begin{center}
\begin{tikzpicture}[>=Stealth,every node/.style={draw,rounded corners,fill=blue!8,minimum width=1.45cm,minimum height=0.55cm,font=\small},edge/.style={->,thick}]
\node (project) at (0,1.4) {project};
\node (plan) at (-2,0) {plan};
\node (build) at (2,0) {build};
\node (req) at (-3.1,-1.6) {requirements};
\node (design) at (-0.9,-1.6) {design};
\node (tests) at (1.1,-1.6) {tests};
\node (docs) at (3.1,-1.6) {docs};
\draw[edge] (project)--(plan);
\draw[edge] (project)--(build);
\draw[edge] (plan)--(req);
\draw[edge] (plan)--(design);
\draw[edge] (build)--(tests);
\draw[edge] (build)--(docs);
\draw[edge,red!75!black,dashed,bend right=22] (req) to node[draw=none,fill=none,below,font=\scriptsize] {cross-link} (tests);
\end{tikzpicture}
\end{center}
\predict{Why is \texttt{requirements -> tests} drawn across the two branches?}
\note{Draw attention to the red dashed cross-link. The six nearby edges show a natural hierarchy, but requirements to tests is a separate dependency. Have students name the seven edges: two from project, two from plan, two from build, and the cross-link.}
\end{frame}

\begin{frame}{Read an edge before you read a trace}
\princ{EDGE MEANING}
\begin{columns}[T]
\begin{column}{0.47\textwidth}
\begin{block}{Dependency question}
Can $B$ be completed before $A$?\\[3pt]
For \texttt{A -> B}, the answer is no.
\end{block}
\end{column}
\begin{column}{0.47\textwidth}
\begin{block}{Containment question}
Is $B$ shown under $A$?\\[3pt]
This is a different claim.
\end{block}
\end{column}
\end{columns}
\vspace{8pt}
\see{one drawing can answer only the question its edges were defined to answer.}
\predict{Does the graph force tests to be drawn under requirements?}
\note{Make the distinction explicit before converting the graph. The dependency graph says requirements must precede tests, but it does not say where a directory tree must draw tests. Ask the prediction; the answer is no. A hierarchy can choose a convenient display without preserving every dependency edge.}
\end{frame}
